\documentclass[11pt,a4paper]{article}

\usepackage[T1]{fontenc}
\usepackage[utf8]{inputenc}
\usepackage[french]{babel}
\usepackage{lmodern}
\usepackage{microtype}
\usepackage[a4paper,margin=2.4cm]{geometry}
\usepackage{enumitem}
\usepackage[autostyle]{csquotes}
\usepackage[hidelinks]{hyperref}
\usepackage{amsmath}
\usepackage[backend=biber,style=authoryear,maxcitenames=2,maxbibnames=12]{biblatex}
\addbibresource{../bibliographie/references.bib}

\setlength{\parindent}{0pt}
\setlength{\parskip}{0.65em}
\setlist[itemize]{leftmargin=1.5em,itemsep=0.35em,topsep=0.35em}
\setlist[enumerate]{leftmargin=1.8em,itemsep=0.5em,topsep=0.4em}

\newcommand{\projectversion}{3}
\newcommand{\projectdate}{28 juillet 2026}

\title{L'émergence de l'intéressant\\
\large Naissance des idées, compréhension et singularité des formes}
\author{François Pachet}
\date{Projet de thèse -- version \projectversion\\\small \projectdate}

\begin{document}

\maketitle

\begin{abstract}
Ce projet vise à constituer l'intéressant comme objet philosophique. Il part
d'une expérience ordinaire mais difficile à stabiliser : une idée surgit, une
forme devient intelligible, un passage retient l'attention et appelle une
reprise. L'hypothèse centrale est qu'est intéressant ce qui, pour un sujet situé
et relativement à un horizon donné, déclenche et soutient une construction
produisant une prise nouvelle sur un objet, un problème ou un espace de
possibilités. L'intéressant n'est donc ni une qualité intrinsèque de la forme,
ni une préférence privée, mais une relation dispositionnelle, historique et
temporelle. La thèse en recherchera la généalogie, en dégagera les propriétés,
puis mettra le concept à l'épreuve dans cinq laboratoires : la musique,
l'intelligence artificielle, la créativité, la philosophie et la littérature.
Sa dimension réflexive sera décisive : le concept de l'intéressant peut
lui-même devenir intéressant et transformer le sujet qui tente de le saisir.
\end{abstract}

\section{Origine, question et objectif}

Ce travail est issu d'une pratique constructiviste et empirique de la recherche
en intelligence artificielle. J'ai passé quarante ans à construire des systèmes
pour générer, transformer ou accompagner des musiques et des textes. Le projet
doctoral ne consiste pas à résumer ces travaux ni à leur donner après coup une
unité philosophique artificielle. Il part d'une question plus exigeante :
\emph{comprendre ce que je fais} lorsque je construis un système, reconnais une
solution, élimine une production correcte mais sans intérêt, ou m'arrête devant
une forme qui semble s'imposer.

La question directrice est la suivante :

\begin{quote}
Sous quelles conditions quelque chose apparaît-il à un sujet comme
intéressant, soutient-il son activité, puis perd-il ou transforme-t-il ce
pouvoir ?
\end{quote}

Une chanson, une démonstration, une phrase ou une idée peut retenir l'attention
de manière disproportionnée, provoquer des comparaisons, des reprises et des
variations. Cet effet ne se réduit ni à la nouveauté brute, puisque l'aléatoire
surprend sans nécessairement intéresser, ni à la beauté, puisqu'une forme
familière peut demeurer belle tout en étant devenue prévisible, ni à la
popularité, ni à une préférence subjective entendue comme donnée sans histoire.

L'enjeu est triple. Il faut d'abord établir la place philosophique de
l'intéressant parmi des notions concurrentes. Il faut ensuite construire un
concept suffisamment déterminé pour produire des distinctions et des
contre-exemples. Il faut enfin montrer ce que ce concept permet de comprendre
dans des pratiques effectives, sans convertir automatiquement un résultat
scientifique ou une expérience personnelle en preuve philosophique.

\part{Constituer l'intéressant comme problème philosophique}

\section{Un angle mort plutôt qu'une absence}

La philosophie n'a pas simplement ignoré l'intéressant. Elle a traité sous
d'autres noms plusieurs de ses dimensions : accès à l'intelligible chez Platon
et Aristote
\parencite{platon2016republique,platon1995theetete,aristote1991metaphysique,aristote2005seconds};
intuition et évidence chez Descartes et Spinoza
\parencite{descartes2002regles,descartes1992meditations,spinoza1966ethique};
jugement chez Kant \parencite{kant2001raison,kant1993juger}; invention et
reconfiguration des problèmes chez Peirce, Poincaré, Hadamard, Bachelard et
Valéry
\parencite{peirce1955writings,poincare1908science,hadamard1959invention,bachelard1999formation,valery1973cahiers};
mémoire, durée et constitution temporelle du sens chez Bergson, Husserl et
Merleau-Ponty
\parencite{bergson2013evolution,bergson2013pensee,husserl1970experience,husserl1957logique,merleau1945phenomenologie};
changement d'aspect chez Wittgenstein \parencite{wittgenstein2004recherches};
symbolisation, différence et répétition chez Goodman et Deleuze
\parencite{goodman1976languages,deleuze1968difference}.

Le diagnostic proposé n'est donc pas celui d'une absence lexicale, mais d'une
\emph{dispersion fonctionnelle}. Le vrai évalue la validité d'une proposition ;
le beau un certain régime du jugement esthétique ; le plaisir une tonalité de
l'expérience ; la surprise un écart à l'attente ; la curiosité une orientation
vers une information manquante ; la pertinence un rapport à une question ou à
un but ; la préférence un ordre de choix. Chacun de ces concepts saisit une
dimension réelle. Aucun n'explique à lui seul pourquoi une forme correcte,
intelligible, vraie ou agréable déclenche chez un sujet une activité qui
transforme ses représentations et ses possibilités d'action.

L'intéressant pourrait nommer ce niveau de sélection et de transformation. Il ne
remplace pas les notions voisines : il permet d'examiner leurs articulations et
leurs divergences. On peut aimer un morceau sans le trouver intéressant, être
surpris sans acquérir de prise nouvelle, reconnaître une proposition vraie sans
juger utile de la poursuivre, ou être mobilisé par un objet pénible et
moralement condamnable.

\section{État de l'art et programme de lecture}

Une généalogie historique est en cours de constitution. Elle devra examiner
notamment l'esthétique de l'intéressant chez Christian Garve, son opposition au
beau et son rapport à la modernité chez Friedrich Schlegel, l'essai
\emph{Ueber das Interessante} de Schopenhauer, le rapport entre ennui,
variation et \emph{det Interessante} chez Kierkegaard, ainsi que l'articulation
de l'intéressant et du vrai chez Whitehead. Le lexique historique de
\textcite{nannini2018interesting} fournit un premier guide, mais il ne dispense
pas du retour aux textes primaires.

Au XX\textsuperscript{e} et au XXI\textsuperscript{e} siècles, quelques travaux
ont directement tenté de conceptualiser l'intéressant. Stace, Kolnai et Ngai en
ont analysé la valeur esthétique, la faiblesse apparente et la circulation
discursive
\parencite{stace1944interestingness,kolnai1964concept,ngai2008merely,ngai2012categories}.
\textcite{epstein2009interesting} insiste sur sa structure transdisciplinaire et
\textcite{grimm2011interesting} sur sa valeur épistémique pour la sélection des
questions. Ces textes empêchent de revendiquer trop vite la découverte d'un
continent vierge. Ils confirment plutôt une faible canonisation : l'intéressant
est présent, mais sans architectonique commune comparable à celles du vrai ou du
beau.

Cet état de l'art restera méthodologiquement stratifié. Une notice
bibliographique vérifiée, un passage repéré, un texte effectivement lu et une
interprétation intégrée à l'argumentation ne constituent pas le même degré
d'acquisition. Les sources déjà lues peuvent soutenir le texte ; les autres
forment un programme explicite de lecture et de vérification. Cette distinction
évite de transformer une suggestion bibliographique en autorité acquise.

\section{L'intéressant et le vrai}

L'intéressant ne saurait devenir un substitut relativiste au vrai. Un souvenir
de discussion avec Heinz Wismann formule brutalement la difficulté :

\begin{quote}
Deleuze et consorts étaient stimulants mais faux. Mais que la philosophie se
devait d'être vraie meme si ennuyeuse
\end{quote}

La citation dissocie deux valeurs. L'intéressant ou le stimulant recommande de
poursuivre une construction ou une enquête ; le vrai autorise à retenir une
proposition comme valide. Une proposition fausse peut être heuristiquement
féconde sans que cette fécondité la rende vraie. Réciproquement, une proposition
vraie peut être ennuyeuse sans cesser d'être vraie.

Cette dissociation ne rend pas l'intéressant philosophiquement superflu. Parmi
les vérités possibles, beaucoup sont triviales ou sans portée pour une enquête
donnée. La vérité ne sélectionne pas à elle seule les problèmes auxquels
consacrer de l'attention. L'intéressant intervient prospectivement dans le choix
des questions, des hypothèses et des constructions ; la vérité intervient dans
leur validation. La philosophie peut partir d'une erreur stimulante, mais elle
ne peut convertir la fécondité de cette erreur en autorisation de demeurer dans
le faux.

\section{L'hostilité de la formalisation}

La science dispose d'instruments puissants pour rendre un problème traitable :
définition opératoire, décomposition, contrôle des variables, répétabilité,
normalisation, moyenne, benchmark et optimisation. Mais ces instruments tendent
aussi à remplacer silencieusement le problème initial par un substitut plus
stable. « Produire quelque chose d'intéressant » devient alors produire quelque
chose de syntaxiquement correct, de sémantiquement cohérent, de probable dans un
corpus, de surprenant ou de conforme à des contraintes.

Il existe en ce sens une \emph{hostilité structurelle de la formalisation à
l'égard de l'intéressant}. Il ne s'agit ni d'une intention hostile des
chercheurs ni d'un argument antiscientifique. La difficulté tient à la nature
relationnelle, historique et transformatrice du phénomène : le rendre mesurable
peut neutraliser la relation, l'histoire et la transformation qui le
constituent. L'avertissement contre la mathématisation prématurée
\parencite{russell1995rationality,russell1995awardlecture} est ici un cas
directeur.

Chaque mesure devra donc être traitée comme une hypothèse locale. Une
formalisation réussie ne doit pas seulement rendre le problème soluble ; elle
doit expliciter ce qu'elle préserve de la question et ce qu'elle abandonne. La
thèse recherchera ainsi une formalisation réflexive, capable de conserver ses
résidus au lieu de les déclarer sans importance.

\part{Construire le concept}

\section{Définition relationnelle et constructive}

Le candidat actuel pour le cœur de la thèse est le suivant :

\begin{quote}
Est intéressant ce qui, pour un sujet situé et relativement à un horizon donné,
déclenche et soutient une construction produisant une prise nouvelle sur un
objet, un problème ou un espace de possibilités.
\end{quote}

La relation minimale peut être notée
\[
  I(F,S\mid H,t),
\]
où \(F\) désigne la forme rencontrée, \(S\) le sujet, \(H\) l'horizon social,
culturel et technique, et \(t\) le moment de la rencontre. L'horizon collectif
ne s'ajoute pas extérieurement aux deux pôles : il forme les attentes et les
compétences du sujet autant que le répertoire des formes disponibles. Il devient
un troisième terme autonome lorsque l'enquête porte sur la circulation ou la
stabilisation sociale de l'intérêt.

Le terme de \emph{construction} désigne ici une activité, parfois implicite, et
non nécessairement la fabrication d'un objet matériel. Elle peut être :

\begin{itemize}
  \item perceptive et conceptuelle, lorsque le sujet discrimine des invariants,
        forme une catégorie ou modifie un découpage ;
  \item explicative, lorsqu'il agence des relations en modèle ou reconstruit le
        problème dont l'objet apparaît comme une solution ;
  \item opératoire, lorsqu'il produit une comparaison, une prédiction, une
        intervention, un algorithme ou un artefact.
\end{itemize}

Une construction arbitraire ne suffit pas. La prise nouvelle doit pouvoir se
manifester par une distinction, un transfert, une anticipation, une variation
contrôlée, une intervention ou un artefact dont les réussites et les échecs sont
examinables. L'intéressant ne garantit pourtant ni la vérité de cette prise ni
sa valeur morale : il en décrit la puissance d'amorçage et de transformation.

\section{Propriétés candidates}

La définition constructive ne constitue pas encore une axiomatique. Elle dégage
un faisceau de propriétés candidates, ni toutes indépendantes ni nécessaires
dans chaque cas. Leur fonction est de produire des variations, des objections
et des comparaisons.

\subsection{Statut relationnel et dispositionnel}

\begin{enumerate}
  \item \textbf{Indexicalité relationnelle.} L'intéressant n'est ni contenu
        dans la forme seule ni réductible à un état privé. Il varie avec
        \(F\), \(S\), \(H\) et \(t\).
  \item \textbf{Dispositionnalité.} Une forme peut être susceptible de
        déclencher une construction sans le faire effectivement, faute
        d'attention, de temps ou de compétences. La disposition ne s'actualise
        que dans certaines conditions de rencontre.
  \item \textbf{Dépendance au chemin.} L'ordre des expériences compte. Une
        forme rencontrée trop tôt peut rester opaque ; après l'acquisition
        d'une nouvelle prise, elle peut devenir intéressante.
\end{enumerate}

\subsection{Dynamique de la construction}

\begin{enumerate}[resume]
  \item \textbf{Écart assimilable.} Il faut assez de différence pour appeler
        une construction et assez de prise pour qu'elle puisse commencer. Le
        familier tend vers la trivialité ; l'étranger absolu vers le bruit ou
        l'abandon.
  \item \textbf{Productivité représentationnelle.} Une forme intéressante
        produit plus de comparaisons, d'hypothèses, de questions, de scénarios
        ou de distinctions qu'elle n'en énonce explicitement.
  \item \textbf{Incomplétude structurée.} Une lacune n'intéresse que si elle
        appelle une continuation orientée. L'objet fournit des indices pour
        construire ce qui manque ou reformuler la question.
  \item \textbf{Transformation de l'espace des possibles.} La construction
        forte ne se contente pas d'ajouter une représentation : elle modifie
        les questions formulables, les distinctions disponibles ou les actions
        concevables.
  \item \textbf{Épluchage constructif.} Le déclenchement ouvre une succession
        de prises : chacune dissipe une opacité, révèle une couche ou reformule
        le problème, puis sert d'instrument à l'opération suivante. Le processus
        se termine par épuisement de l'objet, épuisement du sujet, abandon devant
        une résistance sans prise, clôture décidée d'une construction encore
        féconde ou relance par une résistance productive.
  \item \textbf{Non-monotonie.} Davantage d'information, de surprise ou
        d'attention n'implique pas davantage d'intérêt. Une explication peut
        relancer la construction ou fermer trop vite l'indétermination qui la
        soutenait.
  \item \textbf{Auto-extinction ou relance.} Certaines relations sont
        consumables : la résolution supprime leur intérêt. D'autres sont
        génératives : chaque prise révèle une résistance ou une question
        nouvelle.
\end{enumerate}

Cette dynamique oppose l'intéressant à la contemplation au sens fort retenu ici.
La distinction ne passe pas entre activité corporelle et immobilité : écouter,
lire ou regarder peuvent engager un travail constructif intense. La contemplation
désigne le régime qui suspend les opérations transformatrices afin de maintenir
l'objet intact ou à distance ; elle interdit ou repousse ainsi la construction.
Dès qu'elle produit une distinction réutilisable, une variation ou une prise
nouvelle, elle a déjà basculé dans le régime de l'intéressant.

Cette vie de l'intéressant rejoint le processus créatif et son problème de la
fin \parencite{pachet2026impossibilite}. Les décisions locales ne contiennent
pas nécessairement leur propre critère d'arrêt. Achever peut donc exiger
d'imposer une limite à une construction qui pourrait continuer. Cette clôture
rétroagit sur le parcours et le constitue comme un tout : la fin de la
construction ne signifie pas nécessairement la mort de l'intérêt ni
l'épuisement de l'objet.

\subsection{Valeur, transfert et circulation}

\begin{enumerate}[resume]
  \item \textbf{Neutralité affective relative.} L'intéressant peut être
        agréable, pénible, inquiétant ou répugnant. Il ne constitue pas en
        lui-même une approbation morale ou esthétique.
  \item \textbf{Résonance interdomaines.} Une construction gagne en portée
        lorsqu'elle devient schème, modèle ou métaphore contrôlée pour d'autres
        domaines.
  \item \textbf{Compression générative.} Certaines formes condensent
        rétrospectivement plusieurs phénomènes tout en ouvrant prospectivement
        de nouvelles dérivations.
  \item \textbf{Contagion médiée.} L'intérêt ne se transmet pas comme une
        propriété attachée à l'objet. On peut toutefois préparer chez autrui les
        conditions de son apparition en installant une attente ou en fournissant
        des prises.
\end{enumerate}

\subsection{Propriété réflexive}

\begin{enumerate}[resume]
  \item \textbf{Auto-application.} Le concept de l'intéressant peut lui-même
        déclencher les comparaisons, distinctions, modèles et révisions qu'il
        décrit. « L'intéressant est intéressant » affirme que le concept peut,
        dans certaines conditions, entrer dans son propre domaine
        d'application.
\end{enumerate}

\section{Temporalité, apprentissage et presque-apprenable}

L'intérêt se déplace avec l'apprentissage. Une régularité d'abord opaque peut
devenir accessible, produire un gain de compréhension, puis s'épuiser lorsqu'elle
est acquise. Les théories du progrès de compression et de la motivation
intrinsèque formalisent une version de cette dynamique
\parencite{schmidhuber1997interesting,schmidhuber2009compression,oudeyer2007intrinsic,oudeyer2007typology}.
Si \(E_t(o,S)\) désigne l'erreur du modèle du sujet \(S\) face à l'objet \(o\),
un progrès local peut être noté
\[
  P_t(o,S)=E_{t-\Delta}(o,S)-E_t(o,S).
\]
Les situations déjà maîtrisées et les situations inapprenables ont toutes deux
un progrès proche de zéro. Les niches de progrès sont celles où une amélioration
est possible.

Cette hypothèse rejoint, sans s'y réduire, la zone intermédiaire entre ennui et
anxiété. À compétences fixes, une difficulté croissante fait apparaître la
succession
\[
  \text{ennui}\longrightarrow\text{intéressant/flow}
  \longrightarrow\text{anxiété}.
\]
Pour une tâche fixe, l'apprentissage peut inverser le parcours. La zone de flow
désigne ici un ajustement structurel entre compétence et difficulté, non le seul
état psychologique d'absorption. Une hésitation ou une rupture peuvent rester
intéressantes tant que le sujet conserve une prise et peut réorganiser ses
attentes.

Le langage déplace cette frontière. Un mot, un concept ou une notation compacte
des distinctions et parfois des opérations en une unité réutilisable. Ce qui
exigeait un travail devient une prise pour une construction nouvelle. La zone
du presque-apprenable est donc historiquement mobile et symboliquement médiée
\parencite{vygotsky1978mind,lenat1991thresholds,lenat2008turtle,koestler1989act,steels2005coordinating,steels2015talkingheads}.
Nommer ne suffit toutefois pas : un terme n'est une prise que s'il peut être
déplié en distinctions, comparaisons, prédictions, opérations ou exemples
contrôlables.

\section{Récursivité : « l'intéressant est intéressant »}

La relation possède d'abord une récursivité diachronique :
\[
  I(F,S\mid H,t)\longrightarrow
  \text{construction}\longrightarrow S'
  \longrightarrow I(F,S'\mid H,t+1).
\]
La construction modifie le sujet et donc les conditions des évaluations
suivantes. L'objet n'est plus rencontré par le même sujet au sens pertinent :
ses catégories, ses attentes et son espace de possibilités ont changé.

Elle possède aussi une récursivité métaconceptuelle. Il peut devenir intéressant
de comprendre pourquoi une chose est intéressante, puis de comprendre ce que
cette première explication laisse encore inexpliqué. L'enquête prend alors pour
nouvel objet la relation qui l'avait déclenchée. La formule
\emph{l'intéressant est intéressant} n'est donc pas une tautologie universelle,
mais l'énoncé d'une auto-application possible.

Cette boucle explique partiellement l'insaisissabilité du concept. Toute prise
acquise transforme le répertoire d'exemples, les capacités de discrimination et
la frontière du problème ; une définition féconde fait ainsi apparaître les cas
qui obligeront à la réviser. Il n'en résulte ni paradoxe ni ineffabilité. Des
propriétés locales et contrôlables restent formulables, mais leur théorie doit
être réflexive, historique et révisable.

Enfin, l'auto-application n'est pas un critère de vérité. Une théorie
intéressante de l'intéressant peut être fausse ; une théorie ennuyeuse peut être
exacte. La récursivité décrit ce que l'enquête fait à son objet et à son sujet,
non ce qui valide ses propositions.

\section{Cadre théorique et programme de naturalisation}

Le cadre théorique articulera trois exigences. La première vient d'une
philosophie du sujet : l'expérience dépend de conditions de réception et de
jugement sans être réductible à un trait objectif. La deuxième vient d'une
phénoménologie de la temporalité : rétention, protention, mémoire corporelle et
changement d'aspect décrivent la constitution temporelle d'une forme. La
troisième vient des théories morphodynamiques et des systèmes complexes : une
organisation globale possède des propriétés qui ne se lisent pas dans ses
éléments isolés
\parencite{petitot1993phenomenologie,petitot2002naturaliser}.

La préface de la troisième partie de l'\emph{Éthique} fournit un précédent au
programme de naturalisation \parencite{spinoza1966ethique}. Spinoza refuse de
traiter les affects humains comme un « empire dans un empire » soustrait aux
lois communes de la nature. Étudier l'intéressant reprend ce geste : une
expérience subjective et située peut être expliquée par ses causes sans être
réduite à une propriété intrinsèque de l'objet. Il faut distinguer
l'objectivation du phénomène, qui rend ses conditions et ses trajectoires
comparables, de l'objectivisme qui déposerait l'intérêt dans la forme seule.

Les idées des affections enveloppent à la fois « la nature de notre corps et la
nature présente du corps extérieur »
\parencite[III, prop. XXVII, dém.]{spinoza1861oeuvres}. Ce schème permet de
penser une causalité de la rencontre : structure de la forme, histoire et état
du sujet, horizon collectif et circonstances concourent à produire l'intérêt.
Celui-ci peut être décrit comme un régime affectivo-cognitif de transition qui
modifie la puissance du sujet à percevoir, comprendre, anticiper ou agir.

Naturaliser cette relation exige de coordonner descriptions à la première
personne, comportements, apprentissages, structures formelles et modèles
computationnels, sans postuler qu'un seul niveau épuisera les autres. Le plaisir
de découvrir et celui de construire
\parencite{feynman1999pleasure,florman1996existential} peuvent accompagner
l'intéressant sans le définir : une compréhension pénible peut rester
intéressante, tandis qu'une explication plaisante mais sans pouvoir de
distinction peut n'être qu'une impression de clarté.

\part{Mettre le concept au travail}

\section{La musique : temporalité et micro-émotions}

La musique constitue le laboratoire principal de l'enquête parce qu'elle rend
particulièrement visibles les rapports entre mémoire, attente, surprise,
résolution, répétition et épuisement. Meyer puis Narmour ont montré que l'affect
et la signification musicale dépendent d'attentes apprises, réalisées, retardées
ou déjouées
\parencite{meyer1956emotion,narmour1990basic,narmour1992complexity}. Les modèles
informationnels et neuro-symboliques précisent certains de ces processus
\parencite{abdallah2009information,berger1999expectations,gang1999unified}.

Expliquer pourquoi un titre est intéressant ne consiste cependant pas à repérer
une quantité moyenne de surprise. L'enquête portera sur des trajectoires
d'attention à plusieurs échelles : micro-émotions locales, relations entre
mélodie et harmonie, forme globale, histoire des écoutes et activité de
l'auditeur. Les chansons Jotney fourniront un cas directeur : elles peuvent être
entendues comme des solutions au problème de maintenir une mélodie autonome et
chantable tout en faisant voyager l'harmonie sans exhiber le procédé.

\emph{Histoire d'une oreille} fournit un vocabulaire phénoménologique et une
méthode \parencite{pachet2018oreille}. L'\emph{écoute superposée} soumet la
musique à des variations intérieures. L'\emph{optale} désigne le sentiment
qu'une forme est un optimum local parce que ses variantes proches échouent ; la
\emph{pseudoptale} avertit que la familiarité peut simuler cette nécessité. Les
\emph{fausses évidences} et l'\emph{écoute surattentive} permettent d'examiner
les faux positifs produits par l'auditeur lui-même. Le terrain musical mettra
ainsi à l'épreuve la temporalité, la dépendance au chemin, l'écart assimilable,
la non-monotonie et l'auto-extinction de l'intérêt.

La bonne définition de mots croisés décrite par Perec fournit un analogue
linguistique de l'optale \parencite{perec1999motscroises}. Sa solution paraît
évidente une fois trouvée, alors que le problème semblait auparavant insoluble : le
couple indice--réponse apparaît alors comme un optimum local qu'une formulation
plus explicite ou une autre réponse affaiblirait. Cette propriété est cruciale,
voire \emph{cruciverbiale}, parce que le cas rend particulièrement brève et
observable la transformation qui fait passer d'une pluralité de lectures à une
nécessité rétrospective.

\section{L'intelligence artificielle : expliciter et perdre la cible}

Les systèmes d'intelligence artificielle constituent un deuxième laboratoire.
Les travaux sur Flow Machines, les interactions réflexives, la génération
markovienne sous contraintes et l'\emph{interestingness} computationnelle ont
produit des objets, des protocoles, des réussites et des échecs
\parencite{pachet2006clefs,pachet2008machines,pachet2012virtuosity,pachet2021assisted,pachet2026markov}.
Ils rendent explicites des opérations souvent laissées implicites :
échantillonnage, contrainte, évaluation, construction de modèles, sélection et
direction.

L'IA est également le cas réflexif privilégié de l'hostilité de la
formalisation. Un système peut résoudre le problème formel tout en supprimant sa
partie intéressante. Les métriques de nouveauté, de surprise, de diversité, de
préférence ou de popularité constituent des indicateurs locaux ; elles ne
doivent pas être confondues avec la cible relationnelle
\(I(F,S\mid H,t)\). Le succès expérimental pertinent ne sera donc pas seulement
la génération d'objets jugés intéressants, mais la capacité d'un modèle à
prédire ou expliquer des transformations d'attention, d'apprentissage et de
construction.

\section{La créativité : produire, juger et reconstruire un problème}

La créativité permet de distinguer production du nouveau et production de
l'intéressant. Une génération peut être inédite tout en restant arbitraire ;
une solution peut employer des matériaux connus tout en reconfigurant l'espace
des possibles. Certaines formes paraissent nécessaires après leur apparition,
alors qu'elles n'étaient pas déductibles auparavant. Elles rendent visibles
leurs propres contraintes et donnent l'impression d'une simplicité de surface
qui dissimule une genèse improbable
\parencite{pachet2026impossibilite,pachet2026virtuosite}.

L'objet intéressant peut ainsi être rencontré avant que son problème soit
connu. Le sujet reconstruit à rebours les contraintes que la solution satisfait,
les alternatives qu'elle exclut et la question qu'elle rend visible. Cette
reconstruction ne doit pas devenir ad hoc : elle devra permettre des variations,
des prédictions ou des comparaisons qui auraient pu échouer. Le terrain de la
créativité mettra particulièrement à l'épreuve la productivité
représentationnelle, l'incomplétude structurée, la transformation de l'espace
des possibles et la compression générative.

\section{La philosophie : fécondité, vérité et auto-application}

Les textes philosophiques ne serviront pas d'habillage disciplinaire. Ils
contraindront le concept, révéleront ses présupposés et fourniront des
contre-exemples. Le corpus historique permettra d'étudier comment le phénomène a
été distribué entre jugement, esthétique, curiosité, attention, invention,
vérité et ennui. La remarque de Heinz Wismann empêchera notamment de confondre
la capacité d'un texte à mettre la pensée en mouvement avec la validité de ses
propositions.

La philosophie constitue aussi le terrain de l'auto-application. Une théorie de
l'intéressant doit expliquer pourquoi elle déclenche elle-même des distinctions
et des reprises, tout en acceptant que cette fécondité ne la valide pas. Son
évaluation exigera donc deux séries de critères : rendement conceptuel,
production de problèmes et transferts d'une part ; cohérence, exactitude
historique et résistance aux objections d'autre part.

\section{La littérature : incomplétude, style et interprétation}

La littérature fournit un terrain distinct de la musique parce que la
construction y est médiée par le langage, la narration, les voix et les attentes
sémantiques. Elle permettra d'étudier comment une phrase ou un récit fournit
assez de prises pour orienter l'interprétation sans l'épuiser. L'intérêt peut y
naître d'une incomplétude structurée, d'une attente narrative, d'une contrainte
formelle, d'une ruse avec la langue ou d'un déplacement des catégories
ordinaires.

Ce terrain devra également distinguer fécondité et pseudo-profondeur. Une
opacité peut produire des interprétations parce qu'elle est structurée ; elle
peut aussi entretenir seulement la promesse indéfinie qu'une clef viendra.
Variations stylistiques, paraphrases, continuations concurrentes et comparaison
des lectures permettront d'examiner si le texte produit des prises
transférables ou seulement une fascination. La littérature testera ainsi
l'incomplétude structurée, la contagion médiée, la dépendance à l'horizon
culturel et la récursivité interprétative.

\section{Méthode comparative et critères d'épreuve}

Les cinq terrains n'ont pas vocation à illustrer après coup une théorie déjà
achevée. Chacun doit contraindre le concept, révéler une limite ou produire une
distinction. La recherche combinera :

\begin{itemize}
  \item l'analyse conceptuelle de l'intéressant et de ses voisins ;
  \item la reconstruction historique et la lecture directe des sources ;
  \item une phénoménologie située de cas musicaux, littéraires, mathématiques,
        scientifiques et philosophiques ;
  \item des comparaisons contrefactuelles : variations intérieures,
        continuations compétentes, reprises et imitations ratées ;
  \item l'analyse réflexive de systèmes construits et, lorsque cela est
        pertinent, des expériences sur la reprise, l'abandon, l'apprentissage
        et l'évolution des jugements ;
  \item la confrontation avec les modèles de l'attente, de la compression, de
        la curiosité, du flow et de la génération contrainte.
\end{itemize}

Trois niveaux resteront distingués : les résultats scientifiques et les
propriétés effectives des systèmes ; les propositions philosophiques ; les
articulations argumentées expliquant ce que la construction d'un système ou
l'analyse d'un cas rend pensable sans prétendre qu'elle le démontre.

Une hypothèse gagnera en force si elle permet de décrire une trajectoire,
d'identifier les conditions de variation de l'intérêt, de distinguer des
phénomènes auparavant confondus, d'anticiper un changement ou de produire un
contre-exemple. La thèse assumera aussi les résultats négatifs : certaines
propriétés pourront n'être valables que pour un terrain, et aucun modèle unique
ne sera présupposé capable d'épuiser l'objet.

\section{Corpus}

Le corpus principal réunira \emph{Histoire d'une oreille}, les chansons Jotney,
les essais sur l'impossibilité de créer et la virtuosité
\parencite{pachet2018oreille,pachet2026impossibilite,pachet2026virtuosite}, les
travaux sur les systèmes génératifs et l'\emph{interestingness}, ainsi que les
sources philosophiques et scientifiques effectivement mobilisées.

Le corpus sera organisé en trois cercles : sources directement argumentées dans
le manuscrit ; cas de comparaison et d'objection ; programme de lectures encore
à consolider. Le corpus littéraire sera resserré en fonction des propriétés à
tester plutôt qu'étendu selon un principe d'exhaustivité historique. Cette
organisation doit maintenir un lien explicite entre chaque source, la
proposition qu'elle soutient et le statut réel de sa lecture.

\section{Contributions attendues}

La thèse vise six contributions :

\begin{enumerate}
  \item une histoire problématisée de l'intéressant et de sa dispersion entre
        des concepts concurrents ;
  \item une définition relationnelle, dispositionnelle et constructive,
        exprimée par \(I(F,S\mid H,t)\) ;
  \item un faisceau organisé de propriétés permettant de comparer apparition,
        maintien, circulation, épuisement et relance de l'intérêt ;
  \item une théorie réflexive expliquant à la fois l'auto-application du concept
        et une part de son insaisissabilité ;
  \item une série d'épreuves différentielles dans la musique, l'IA, la
        créativité, la philosophie et la littérature ;
  \item une méthode de formalisation locale qui distingue explicitement cible,
        indicateur, résultat scientifique et portée philosophique.
\end{enumerate}

Le résultat ne sera pas nécessairement une définition fermée. Il pourra prendre
la forme d'un système articulé de propriétés, de distinctions et de critères de
variation. Un prolongement formel examinera notamment la
non-compositionnalité, la non-monotonie, les modalités dynamiques et les
conditions d'une axiomatique minimale.

\section{Plan courant de la thèse}

\begin{enumerate}
  \item \textbf{Première partie --- Constituer le problème}
    \begin{enumerate}
      \item L'angle mort et les concepts concurrents.
      \item Généalogie et état de l'art de l'intéressant.
      \item L'intéressant et le vrai.
      \item L'hostilité de la formalisation.
    \end{enumerate}
  \item \textbf{Deuxième partie --- Construire le concept}
    \begin{enumerate}
      \item La relation \(I(F,S\mid H,t)\) et la construction d'une prise.
      \item Les propriétés relationnelles, dynamiques et transmissives.
      \item Apprentissage, flow, presque-apprenable et langage.
      \item Récursivité et auto-application : « l'intéressant est intéressant ».
      \item Naturalisation, limites et modèles locaux.
    \end{enumerate}
  \item \textbf{Troisième partie --- Mettre le concept au travail}
    \begin{enumerate}
      \item Musique : attentes, micro-émotions et variations.
      \item Intelligence artificielle : génération, évaluation et perte de la
            cible.
      \item Créativité : nécessité inventée et problème implicite.
      \item Philosophie : fécondité, vérité et réflexivité.
      \item Littérature : incomplétude, style et interprétation.
      \item Comparaison des terrains et révision du concept.
    \end{enumerate}
\end{enumerate}

\section{Travaux personnels directement liés}

\begin{itemize}
  \item \textcite{pachet2026virtuosite} analyse les formes de virtuosité et leurs
        effets recherchés sur les auditeurs.
  \item \textcite{pachet2026impossibilite} examine les obstacles scientifiques
        et conceptuels à l'idée de créer du nouveau.
  \item \textcite{pachet2026markov} synthétise les travaux sur la modélisation
        markovienne et la génération contrôlée par contraintes.
  \item \textcite{pachet2021assisted} tire des projets Flow Machines de nouvelles
        catégories du nouveau.
  \item \textcite{pachet2018oreille} décrit l'ontogenèse d'une oreille musicale
        et fournit les concepts attentionnels mobilisés ici.
  \item \textcite{pachet2012virtuosity} relie virtuosité, créativité et
        construction du système Virtuoso.
  \item \textcite{pachet2008machines} présente les systèmes réflexifs comme des
        machines favorisant la sortie de soi.
  \item \textcite{pachet2006clefs} propose une première modélisation des mélodies
        intéressantes.
\end{itemize}

\printbibliography[title={Références}]

\end{document}
